\centerline{\bf \Huge ОТА. Делимость II}

\begin{flushright}
        \scriptsize{Выживает тот, у кого есть воля к жизни.\\ Мисато Кацураги}
\end{flushright}

\problem{1}
Перемножили несколько натуральных чисел и получили 224, причём самое маленькое число было ровно вдвое меньше самого большого. Сколько чисел перемножили?

\problem{2}
Существует ли целое число, произведение цифр которого равно  а) 1980?  б) 1990?  в) 2000?

\problem{3}
По кругу расставили 2017 чисел. Может ли быть так, что отношение любых двух последовательных чисел - простое число?

\problem{4}
Натуральное число назовём удивительным, если самый большой его собственный делитель (т.е. делитель, не равный 1 и самому числу) на 1 больше самого маленького собственного делителя. Найдите все удивительные числа.

\problem{5}
Опишите все натуральные числа, у которых нечётное количество делителей.

\problem{6}
На доске записаны двузначные числа. Каждое число составное, но любые два числа взаимно просты. Какое наибольшее количество чисел может быть записано?
